\documentclass[11pt,a4paper]{article}

\usepackage[utf8]{inputenc}
\usepackage[T1]{fontenc}
\usepackage[spanish,es-nodecimaldot]{babel}
\usepackage[a4paper,margin=2.2cm]{geometry}
\usepackage{xcolor}
\usepackage{listings}
\usepackage{enumitem}
\usepackage{booktabs}
\usepackage{hyperref}
\usepackage{tcolorbox}
\usepackage{parskip}

\definecolor{azul}{HTML}{174A7E}
\definecolor{azulclaro}{HTML}{EAF3FA}
\definecolor{verdeclaro}{HTML}{EEF7EE}
\definecolor{grisfondo}{HTML}{F5F6F7}
\definecolor{verde}{HTML}{2E7D32}
\definecolor{naranja}{HTML}{A64B00}

\hypersetup{
  colorlinks=true,
  linkcolor=azul,
  urlcolor=azul,
  pdftitle={Laboratorio IA 7: Red neuronal evolutiva}
}

\lstdefinestyle{python}{
  language=Python,
  basicstyle=\ttfamily\footnotesize,
  backgroundcolor=\color{grisfondo},
  frame=single,
  rulecolor=\color{black!20},
  keepspaces=true,
  showstringspaces=false,
  breaklines=true,
  columns=fullflexible,
  keywordstyle=\color{azul}\bfseries,
  commentstyle=\color{verde},
  stringstyle=\color{naranja}
}

\newtcolorbox{objetivo}{
  colback=azulclaro,
  colframe=azul,
  boxrule=0.8pt,
  arc=2pt,
  title=Objetivo de la sesion,
  fonttitle=\bfseries\color{azul}
}

\newtcolorbox{idea}{
  colback=verdeclaro,
  colframe=verde,
  boxrule=0.8pt,
  arc=2pt,
  title=Idea clave,
  fonttitle=\bfseries\color{verde}
}

\newtcolorbox{reto}[1]{
  colback=white,
  colframe=azul,
  boxrule=0.8pt,
  arc=2pt,
  title=#1,
  fonttitle=\bfseries\color{azul}
}

\title{\textbf{Laboratorio IA 7: Red neuronal pequena y evolucion por algoritmo genetico}}
\author{Materia: Inteligencia Artificial}
\date{}

\begin{document}
\maketitle

\begin{center}
  \textbf{Dinamica:} practica en terminal\quad
  \textbf{Entregable:} experimento, codigo, presentacion de resultados y conclusiones
\end{center}

\begin{objetivo}
Construir una red neuronal pequena con PyTorch que controle al robot de percepcion local y evolucionar sus pesos mediante un algoritmo genetico. El mejor individuo se guardara en disco para probarlo posteriormente sobre un mapa modificado a mano.
\end{objetivo}

\section{Continuidad con las clases anteriores}
El laboratorio conserva el mundo del robot usado en las sesiones anteriores: una cuadricula, un punto de inicio, una meta, obstaculos, sensores locales y una funcion de fitness. La progresion es:

\begin{itemize}[leftmargin=*]
  \item En las primeras clases se evoluciono una secuencia de movimientos.
  \item Despues se incorporaron obstaculos y se compararon mutacion y cruce.
  \item Con la percepcion local, el robot paso a consultar el entorno antes de decidir.
  \item En este laboratorio se conserva esa politica reactiva, pero se reemplaza la tabla de reglas por una red neuronal.
\end{itemize}

\begin{idea}
El algoritmo evolutivo sigue seleccionando comportamientos por fitness. Lo que cambia es la representacion del individuo: antes eran comandos o reglas; ahora son los pesos de una red neuronal.
\end{idea}

\section{Preparacion del entorno}
El programa necesita Python, PyTorch y el archivo \texttt{robot\_red\_neuronal\_evolutiva.py}. Con el interprete del laboratorio se puede comprobar la instalacion:

\begin{lstlisting}[style=python]
python3 -c "import torch; print(torch.__version__)"
\end{lstlisting}

Si PyTorch no esta instalado en ese interprete:

\begin{lstlisting}[style=python]
python3 -m pip install --index-url https://download.pytorch.org/whl/cpu torch
\end{lstlisting}

\section{Mapa y sensores}
El mapa se modifica directamente al principio del archivo Python:

\begin{lstlisting}[style=python]
TAMANO = 8
INICIO = (0, 0)
META = (7, 7)
OBSTACULOS = {
    (0, 3), (1, 3), (2, 0), (2, 2), (2, 3),
    (4, 5), (5, 2), (5, 3), (5, 4), (6, 6),
}
\end{lstlisting}

En cada paso el robot consulta las celdas vecinas en orden \texttt{U D L R}. Una celda libre se representa con 0; un borde o un obstaculo, con 1. Por tanto, la entrada de la red tiene cuatro valores binarios. El robot no observa el mapa completo ni recibe directamente la posicion de la meta.

\section{La red neuronal}
La arquitectura es deliberadamente pequena:

\begin{center}
\begin{tabular}{c|c|c}
\toprule
Entrada & Capa oculta & Salida \\
\midrule
4 sensores & 8 neuronas con ReLU & 4 acciones \\
\texttt{U,D,L,R} & combinacion no lineal & un valor por accion \\
\bottomrule
\end{tabular}
\end{center}

La salida no es una accion directamente. Son cuatro valores llamados logits. La accion elegida es la que tenga el valor maximo:

\begin{lstlisting}[style=python]
logits = red(entrada)[0]
indice = int(torch.argmax(logits).item())
accion = COMANDOS[indice]
\end{lstlisting}

Cada individuo de la poblacion es una instancia distinta de \texttt{RedPercepcion}. Sus pesos funcionan como el material genetico.

\section{Fitness y algoritmo evolutivo}
Cada red recorre el mapa durante como maximo 30 pasos. El fitness utilizado es:

\[
F = 500 - 35d + 4p - 30c - 3r + 2000m
\]

Donde $d$ es la distancia Manhattan a la meta, $p$ los pasos utiles, $c$ los choques o movimientos bloqueados, $r$ las posiciones repetidas y $m$ vale 1 cuando se alcanza la meta y 0 en otro caso.

En cada generacion:

\begin{enumerate}[leftmargin=*]
  \item Se evalua toda la poblacion.
  \item Se ordenan los individuos por fitness.
  \item Se conservan seis elites.
  \item Se seleccionan padres entre los mejores individuos.
  \item Se aplica cruce en un punto y mutacion gaussiana sobre los pesos.
  \item Se completa la siguiente poblacion hasta llegar a 40 individuos.
\end{enumerate}

Por defecto se ejecutan 40 generaciones y la semilla es aleatoria. Se puede fijar una semilla para reproducir un experimento concreto.

\section{Ejecutar el laboratorio}
Ejecucion normal, con 40 generaciones, semilla aleatoria y animacion:

\begin{lstlisting}[style=python]
python3 robot_red_neuronal_evolutiva.py
\end{lstlisting}

Ejecucion reproducible y sin pausa entre cuadros:

\begin{lstlisting}[style=python]
python3 robot_red_neuronal_evolutiva.py \\
  --semilla 7 --generaciones 40 --sin-pausa
\end{lstlisting}

La animacion del mejor individuo aparece despues de la evolucion. Cada cuadro muestra la observacion binaria, la accion elegida y la posicion actual.

\section{Guardar y cargar el mejor individuo}
Para entrenar y guardar la red:

\begin{lstlisting}[style=python]
python3 robot_red_neuronal_evolutiva.py \\
  --generaciones 40 --guardar-mejor mejor_red.pt
\end{lstlisting}

El archivo \texttt{mejor\_red.pt} contiene los pesos de PyTorch y los metadatos basicos de la red. Para cargarlo sin evolucionar otra vez:

\begin{lstlisting}[style=python]
python3 robot_red_neuronal_evolutiva.py \\
  --cargar-red mejor_red.pt
\end{lstlisting}

\begin{idea}
Guardar la red separa dos fases del trabajo: primero se aprende una politica y despues se prueba esa misma politica bajo nuevas condiciones.
\end{idea}

\section{Probar en un mapa distinto}
\begin{enumerate}[leftmargin=*]
  \item Ejecuta una evolucion y guarda \texttt{mejor\_red.pt}.
  \item Cambia manualmente \texttt{OBSTACULOS}, \texttt{INICIO} o \texttt{META} en el archivo.
  \item Conserva el mismo \texttt{TAMANO} y evita colocar inicio o meta sobre un obstaculo.
  \item Ejecuta con \texttt{--cargar-red mejor\_red.pt}.
  \item Registra si la red llega, cuantos pasos usa, cuantos choques produce y si entra en un ciclo.
\end{enumerate}

Para que la comparacion sea valida, no se deben cambiar los pesos entre las dos pruebas. Solo se modifica el mapa.

\begin{reto}{Preguntas para discutir}
\begin{enumerate}[leftmargin=*]
  \item Que se conserva entre la clase de percepcion local y este laboratorio?
  \item Que informacion representa un peso de la red que antes representaba una regla?
  \item Por que una red entrenada en un mapa puede fallar en otro mapa?
  \item La politica aprendida generaliza o solo memoriza respuestas locales?
  \item Que cambios permitirian que el robot conociera la direccion de la meta?
\end{enumerate}
\end{reto}

\section{Experimento y tabla de resultados}
Realiza tres entrenamientos con semillas distintas. Guarda cada mejor individuo con un nombre diferente y prueba cada uno en el mapa original y en un mapa modificado.

\begin{center}
\begin{tabular}{c|c|c|c|c|c|c}
\toprule
Semilla & Mapa & Llego & Generacion & Fitness & Choques & Pasos utiles \\
\midrule
7 & Original & & & & & \\
7 & Modificado & & & & & \\
21 & Original & & & & & \\
21 & Modificado & & & & & \\
42 & Original & & & & & \\
42 & Modificado & & & & & \\
\bottomrule
\end{tabular}
\end{center}

\section{Entrega}
Entrega los siguientes elementos:

\begin{itemize}[leftmargin=*]
  \item \texttt{robot\_red\_neuronal\_evolutiva.py}.
  \item Un archivo \texttt{.pt} con al menos un mejor individuo.
  \item La tabla de resultados completa.
  \item Una captura o descripcion de la trayectoria en el mapa original y en el modificado.
  \item Una conclusion de 8 a 10 lineas sobre evolucion, percepcion local y generalizacion.
\end{itemize}

\section{Codigo completo}
\lstinputlisting[style=python]{robot_red_neuronal_evolutiva.py}

\end{document}
